\documentclass[11pt, parskip=full, DIV=14, headsepline]{scrreprt}

\usepackage{jdcommands}
\jdfontsetup{}

\usepackage{amsmath, amssymb}
\usepackage[USenglish]{babel}

\usepackage{pgfplots}
\pgfplotsset{compat=newest}
\usepgfplotslibrary{dateplot}
\usepackage{pgfcalendar}

\usepackage{multicol}
\usepackage{wrapfig}

\usepackage{xparse}

\newlength{\oldintextsep}
\setlength{\oldintextsep}{\intextsep}

\setlength\intextsep{0pt}

\NewDocumentEnvironment{mywraptable}{O{}mO{\wrapoverhang}mO{0pt}}{%
  \setlength{\intextsep}{#5}%  Defaults to 0pt \plus 0pt \minus 0pt
  \wraptable[#1]{#2}[#3]{#4}%
}{%
  \endwraptable%
  \setlength{\intextsep}{\oldintextsep}%  Defaults to 0pt \plus 0pt \minus 0pt
}

\let\oldcenter\center
\let\oldendcenter\endcenter
\renewenvironment{center}{\setlength\topsep{6pt}\oldcenter}{\oldendcenter}

\usepackage{subcaption}
\usepackage{booktabs}
\usepackage{float}
\usepackage{listings}
\lstset{numbers=left, numberstyle=\tiny, frame=lines, breaklines, caption=\lstname}

\usetikzlibrary{external}
\tikzexternalize[prefix=tikz/]

\definecolor{dbblue}{RGB}{0,24,168}
\definecolor{lightblue}{RGB}{0,152,219}
\definecolor{warmgrey}{RGB}{152,143,134}
\definecolor{orange}{RGB}{255,161,0}
\definecolor{green}{RGB}{122,184,0}
\definecolor{palegrey}{RGB}{224,230,230}
\definecolor{midgrey}{RGB}{137,150,160}

\renewcommand{\familydefault}{\sfdefault}
\setlength{\textfloatsep}{8pt plus 1.0pt minus 2.0pt}

\usepackage{scrlayer-scrpage}
\automark[section]{chapter}
%\renewcommand{\chapterpagestyle}{scrheadings}
\pagestyle{scrheadings}
\chead{\headmark}
\cfoot[\pagemark]{\pagemark}

\newcommand*{\FSfont}[1]{}
\newcommand*{\plogo}{\includegraphics[width=6cm]{jdlogo.jpg}}
\newcommand*{\rotrt}[1]{\rotatebox{90}{#1}}
\newcommand*{\rotlft}[1]{\rotatebox{-90}{#1}}
\newcommand*{\topb}{%
  \resizebox*{\unitlength}{\baselineskip}{\rotrt{$\}$}}}
\newcommand*{\botb}{%
  \resizebox*{\unitlength}{\baselineskip}{\rotlft{$\}$}}}
\newcommand*{\titleBC}{\begingroup
\thispagestyle{scrplain}
\FSfont{5jr}% Fontsite Jenson Recut (Centaur)
\begin{center}
\def\CP{\textit{\Huge Active Fund Solutions for a safe future}}
\settowidth{\unitlength}{\CP}
{\color{black}\topb} \\[\baselineskip]
\textcolor{lightblue}{\CP} \\[\baselineskip]
%{\color{black}\LARGE JD Investments} \\
\plogo\\[0.5\baselineskip]
{\color{black}\botb}
\end{center}
\vfill
\begin{center}
{\LARGE\textbf{Andrej Bajic, Nikolay Balantsev, Julian Debus, Clemens Reeh, Nicolay Schmitt, Sabrina Taneja}}\\
\vfill
\today
\end{center}
\endgroup}

\usepackage[skins,many]{tcolorbox}
\newtcolorbox{mybox}[1]{
colframe=midgrey,
colback=palegrey,
coltitle=white,
fonttitle=\bfseries,
width=\textwidth,
left=0mm,
sharp corners,
adjusted title=#1
}

\title{Portfolios for a safe future}
\author{JD POWER}
\begin{document}
\titleBC

\pgfplotsset{y label style={at={(axis description cs:-0.08,.5)}, text depth=0.25ex, anchor=south,}}

\pgfplotsset{
  /pgfplots/bar cycle list/.style={/pgfplots/cycle list={%
      {fill=lightblue, lightblue!80!black},
      {fill=warmgrey, warmgrey!80!black},
      {fill=orange, orange!80!black},
      {fill=green, green!80!black},
      {fill=palegrey, palegrey!80!black},
  }
},
}

\pgfplotsset{
  gen/.style={
      no marks,
      date coordinates in=x,
      xticklabel=\year,
      xtick={2008-01-01,2009-01-01,2010-01-01,2011-01-01,2012-01-01, 2013-01-01, 2014-01-01, 2015-01-01},
      yticklabel style={/pgf/number format/fixed,/pgf/number format/precision=2},
      every axis plot/.append style={thick, smooth},
      legend columns=5,
      every axis legend/.append style={at={(0.5,1.03)}, anchor=south},
      trim axis left,trim axis right,
      stack plots=y,
      cycle list name=exotic,
      area style,
      every axis/.append style={font=\footnotesize},
      legend style={font=\footnotesize},
      legend pos=north west,
  },
  total/.style={red, ultra thick, line legend},
   relw/.style 2 args={
      gen,
      ylabel=,
      yticklabel=\pgfmathparse{100*\tick}\pgfmathprintnumber{\pgfmathresult}\,\%,
      ymin=#1, ymax=#2,
      width=\textwidth,
      height=0.5\textwidth,
      %every axis legend/.append style={
        %at={(0.5,1.03)},
      %anchor=south},
    },
   w/.style 2 args={
      gen,
      ylabel=Performance,
      yticklabel=\pgfmathparse{\tick}\pgfmathprintnumber{\pgfmathresult}\,\%,
      ymin=#1, ymax=#2,
      legend columns=1,
      legend pos=outer north east,
      width=0.8\textwidth,
      height=0.35\textwidth,
      %legend pos=outher north east,
    },
  bottom/.style={axis y line = none},
  cycle list name=exotic,
  assetplot/.style={
    no marks,
    date coordinates in = x,
    xticklabel = \year,
    xtick = {2009-01-01,2011-01-01,2013-01-01,2015-01-01},
    yticklabel style = {/pgf/number format/fixed, /pgf/number format/precision=3},
    yticklabel = \pgfmathparse{\tick}\pgfmathprintnumber{\pgfmathresult}\,\%,
    width = 0.5\textwidth,
    height = 0.3\textwidth,
    trim axis left,
    trim axis right,
    every axis plot/.append style={thick, smooth, dbblue},
  }
}

\chapter{Become active regarding your wealth}
First we would like to introduce JD Investments and ourselves.
JD Investments is one of the leading investment organizations.
The company offers institutions traditional and alternative investments across all major asset classes.
We offer our clients first-rate product management and a consistently superior client service platform.
 
We are JD Investments’ team of most experienced portfolio managers.
With a high level of market experience, global expertise and one of the smallest but most successful asset management team for active asset allocation and risk mitigation in Europe.
The following factors make us specialists in our industry: 

\begin{itemize}
\item[$\checkmark$] Diversified active portfolio manager
\item[$\checkmark$] Global investment and research capabilities 
\item[$\checkmark$] Consultative local delivery 
\item[$\checkmark$] Forward looking investment strategies 
\end{itemize}

We would like to guide you through our investment process and advise you with an investment that is tailored to your individual needs.
You came to us as an experienced investor, looking to expand your investment portfolio with the following five assets: DAX, Dow Jones Industrial Average, Nikkei 225, iPath S\&P 500 VIX Short-Term Futures ETN and the risk free rate (Euribor).
We have compiled an overview of various portfolio strategies that would fit to your risk mentality.
Furthermore, we backtested these strategies, in order to see which of these would fit best to your return expectations.
  
In putting together our portfolio proposals, we worked with a special focus on meeting your demands that are summarized as follows:
Portfolios should (a) have a low turnover/transaction costs, (b) not allow for big drawdowns, (c) maximal Sharpe Ratio under these constraints.
Accordingly, we have set up constraints to (a) limit turnover to $10 \pm 2\%$, p.a., (b) put a cap on maximum drawdown of less than 10\% p.a., (c) maximize realized Sharpe ratio.

In the following proposal we will start off by introducing the asset classes you have selected and our view on them.
The proposal includes a detailed description of our investment philosophy, since it is essential to understand the optimization procedure and the tools used for portfolio construction.
This is followed by a summary of selected, tailor-made portfolios for you, as well as our personal recommendation.


\chapter{Asset Universe}

\tikzsetnextfilename{assets_dax}
\begin{wrapfigure}{r}{0.5\textwidth}
  \centering
  \begin{tikzpicture}
    \begin{axis}[assetplot]
      \addplot[dbblue] table[x=time, y=value.DAX] {plot-data/assets_vxx.txt};
    \end{axis}
  \end{tikzpicture}
\end{wrapfigure}
%
\paragraph{DAX} is a German capitalization-weighted (in EUR) stock market index that includes the top 30 (blue chip) German companies trading on the Frankfurt Stock Exchange.
The DAX is considered to be one of the most important indices for European investors, since the German economy is the fourth largest in the world/the largest in Europe.
Furthermore, many DAX companies are heavily export oriented, giving investors broad exposure to global markets. In fact, more than 75\% of DAX revenues stem from international markets, indicating a sound international diversification.

\tikzsetnextfilename{assets_dj}
\begin{wrapfigure}{r}{0.5\textwidth}
  \centering
  \begin{tikzpicture}
    \begin{axis}[assetplot]
      \addplot[dbblue] table[x=time, y=value.Dow.Jones] {plot-data/assets_vxx.txt};
    \end{axis}
  \end{tikzpicture}
\end{wrapfigure}
%
\paragraph{Dow Jones Industrial Average} is a price-weighted (in USD) stock market index that describes the stock performance of 30 large publicly owned companies based in the United States.
The index proved itself as very stable and able to overcome large market downturns, since most companies in the index are well-established and internationally diversified.
Given America's status as the world's leading economy, the Dow Jones offers investors the chance to participate in the growth of the most influential companies in the US.

\tikzsetnextfilename{assets_nikkei}
\begin{wrapfigure}{r}{0.5\textwidth}
  \centering
  \begin{tikzpicture}
    \begin{axis}[assetplot]
      \addplot[dbblue] table[x=time, y=value.Nikkei] {plot-data/assets_vxx.txt};
    \end{axis}
  \end{tikzpicture}
\end{wrapfigure}
%
\paragraph{Nikkei 225} is a price-weighted (in Yen) stock market index for the Tokyo Stock Exchange and is considered as the best average representative of the Japanese Market.
The Japanese economy is the world’s third largest and is assumed to enter a recovery phase after two decades of stagnation.
Japan's economy leads all world economies in various measures that are testament for its impressive potential for innovation, such as R\&D as proportion of GDP, number of registered patents and environment for the development of business.
Moreover, Japan is considered as mature market gateway to quickly growing Asian market.

\paragraph{Euribor} is an estimate rate at which European banks borrow funds from one another.
It is calculated as a truncated mean of quoted rates.
As of November 10, the Euribor rate was -0.004\%, averaging 0\% throughout November.

\tikzsetnextfilename{assets_vxx}
\begin{wrapfigure}{r}{0.5\textwidth}
  \centering
  \begin{tikzpicture}
    \begin{axis}[assetplot]
      \addplot[dbblue] table[x=time, y=value.VXX] {plot-data/assets_vxx.txt};
    \end{axis}
  \end{tikzpicture}
\end{wrapfigure}
%
\paragraph{iPath S\&P 500 VIX Short-Term Futures ETN} is made of medium-term uncollateralized notes, created by Barclays. Its performance follows the performance of the underlying VIX futures (S\&P 500 VIX ShortTerm Futures Index TR) with rolling maturities of 1 and 2 months.
The main benefit of VXX is that an investor gets exposed to VIX futures without having to constantly roll them over, thereby saving on transaction cost. The correlation between VXX and VIX is almost perfect.
VIX benchmark index reflects market opinion and corresponding price of volatility of S\&P 500 Index over 30 days in the future.


Normally, in times of uncertainty the value of the index increases, so investor can benefit from going long VIX futures when realized volatility appears to be larger than implied.
Short positions in the VIX are taken when investors expect market implied volatility to decrease in the future.
Moreover, VIX has demonstrated a strong negative correlation with the underlying index.
That is, in 80\% of time VIX and S\&P 500 moved in opposite direction.
Therefore, when markets fall, long position in VIX futures should increase in value.
All this is also true for VXX, but compared to futures strategies (with constant rolling over), here we just take long or short position of VXX in our portfolio and slightly adjust it with time.


\section{Our Benchmark: MSCI}
In order to show that we have a good performance we need to compare our performance with a benchmark.
Since we invest in indices of developed countries, the MSCI world index can be used as a benchmark.
The MSCI world is an index  which includes medium and large capitalization stocks of 23 developed markets.
Concerning the performance of the index in our out-of-sample period starting from 1st January 2013 till 30th June 2015 the annual return of the index was 20.597\%, annual volatility -17.867\% and corresponding Sharpe ratio –1.15.

\chapter{Investment Philosophy \& Investment Approach}
\section{Our Philosophy}
Supporters of passive investing believe in the implications of the Capital Asset Pricing Model and market efficiency, which means that market prices fully reflect all fundamentals and hence an outperformance cannot be achieved.
On the other hand, advocates of active investing believe that one can outperform the market.
In this regard it makes sense to consider timing and selection in the investment process.
JD investments believes that one can generate return based on a careful strategic asset allocation and consequent tactical asset allocation.
We believe that the advantages of active investing outweigh those of passive investing.
A passive investor can unwillingly be very concentrated even though he wants to exploit all diversification benefits.
That is the case if the market cap in an index increases during the time he holds the investment.
With active investing we see in our asset universe a diversification benefit.

An actively managed portfolio may have the chance of outperforming the market since it weights the assets differently and may not necessarily own all of the assets in the benchmark index.
It can also use special techniques such as taking short positions in certain assets, which is not possible in a passive indexing strategy.
We as active portfolio managers want to understand our investor’s investment objectives and risk tolerance.
We use this knowledge to define an investment policy with risk and return goals, select a performance benchmark, choose an asset mix and allow for adjustments during the course of the market or changes in investment attitude of the investor.

By choosing the indices you wish to invest in and defining your risk tolerance as well as constraints you have already made your first step as an active investor.
JD Investments supports you in the next steps of active investing by providing sophisticated systematic solutions tailor made for your views and risk preferences.

 
\section{Investment Approach and Assumptions}
JD Investments uses modern portfolio optimization techniques based on most recent research combined with subjective experience of its experts with proven track records in order to manage its funds.
Our culture of risk management and regular supervision of the portfolios is essential for our value adding client solutions.

We use a \textit{time horizon} of three years of daily data.
This implies a rolling window of historical data at each optimization since only then recent information has a similar impact on optimizations.
Furthermore, it is assumed that the added value of older data might reflect a very different market environment and therefore reduces the validity of information.

To adjust the portofolio’s asset allocation we can provide our investors \textit{rebalancing} solutions to realign the weightings on a monthly or yearly basis.

In order to improve our out-of-sample performance we use selected regularization techniques and constraints.
In the following we give our investors a brief overview of the approaches and their rationale.
 
For portfolios with assets traded in different markets at different times we adjust the variance/covariance estimation based on the procedure of \textit{Hayashi and Yoshida}.
This technique improves our estimation of the diversification benefit of assets being part of our asset universe.

Since past performance is a weak indicator for future returns we adjust our expected returns by \textit{shrinking} them towards the average mean.
Moreover, we use shrinkage as specified by Ledoit and Wolf to adjust the variance/covariance estimators.
Shrinkage leads to a better out-of-sample performance due to an increased stability of major optimization inputs.

In order to incorporate our market views in our systematic approach we use the \textit{Black Littermann Model} for certain portfolios.
Views on future developments as well as volatilities on those views are additionally used in order to obtain optimal estimators.

We are further able to propose our investors \textit{robust portfolio choices} which are based on robust statistics.
Our procedure follows Wang and Raftery’s approach\footnote{N. Wang, A. Raftery, 2002. Nearest neighbour variance estimation (NVVE). Robust covariance estimation via nearest neigbor cleaning (with discussion), Journal of the American Statistical Association, Volume 97, Pages 994-1019.} to find robust variance/covariance estimates.
The advantage of this procedure is that it makes results from our estimators much more robust against incorrect market assumptions.

In order to \textit{reduce the level of turnover} and hence cut transaction costs several techniques are selectively used.
Reducing active exposure by keeping some positions fixed or changing to larger optimization frequencies both leads to a smaller turnover.
Further, maximum allocation constraints are used to drive down turnover and reduce transaction costs in each point in time of the optimization.
Generally, our procedure is based on the method of Goldfarb and Idnani\footnote{D. Goldfarb, A. Idnani, 1983. A numerically stable dual method for solving strictly convex quadratic programs, Mathematical Programming, Volume 27, Pages 1-33.}, enabling us to include both short-sale and maximum allocation constraints in our optimization.

To prevent large drawdowns in certain portfolio strategies, market exposure of the fund is reduced by investing parts of the portfolios in assets with low volatility e.g. the risk free asset.

Further, long positions in implied volatility are used to exploit negative correlation to stocks in order to further reduce large drawdowns.
For risk taking investors we also provide solutions where volatility is sold as one source of return in order to leverage on opportunities.
In general implied volatility as an asset is treated with highest care in JD Investments Funds as it reacts completely different than value assets.
Implied volatility as an asset is solely used in order to increase or decrease volatility of a fund such that it matches the investors demand and risk profile.

We provide our solutions in EUR denomination.
Currency risks are not hedged in our funds since we do not have an active view about them.
For clients requesting investments without currency exposure hedged share classes can be made available.

\vspace{-5em}
\chapter{Portfolios}

\section{JD Investments Mean Variance+ Equity Fund}
\begin{mywraptable}{r}{0pt}
  \begin{tabular}{lr}
  \toprule
  Ann. Return & 11.61\%\\
  Ann. Volatility & 8.35\%\\
  Sharpe & 1.39\\
  Max. Drawdown& 5.88\%\\
  Turnover & 10.34\%\\
  \bottomrule
\end{tabular}
\end{mywraptable}
\textit{JD Investments Mean Variance+ Equity Fund} has long exposure to the American, German and Japanese equity markets as well as implied volatility of the US market.
The asset allocation is aimed at minimizing volatility while not having more than 50\% exposure to a single market and not allowing for short sales.
Returns are shrunk in order to reduce errors in expected return assumptions.
The portfolio is very immune against stressed market conditions, as we also include the risk-free asset into our portfolio.% We further enhance diversification due to long positions in volatility.
%
\begin{center}
\tikzsetnextfilename{pf_10}  
\begin{tikzpicture}
  \begin{axis}[w={0}{150}, y filter/.expression={max(y,0)}, legend columns=1]
    \addplot table[x=time, y=value.DAX] {plot-data/mvo.txt}\closedcycle;
    \addlegendentry{DAX}
    \addplot table[x=time, y=value.Dow.Jones] {plot-data/mvo.txt}\closedcycle;
    \addlegendentry{Dow Jones}
    \addplot table[x=time, y=value.Nikkei] {plot-data/mvo.txt}\closedcycle;
    \addlegendentry{Nikkei}
    \addplot table[x=time, y=value.VXX] {plot-data/mvo.txt}\closedcycle;
    \addlegendentry{VXX}
    \addplot table[x=time, y=value.mixed] {plot-data/mvo.txt}\closedcycle;
    \addlegendentry{Euribor}
    \addlegendimage{total} \addlegendentry{Total Value}
  \end{axis}
  \begin{axis}[w={0}{150}]
    \addplot[total] table[x=time, y=value] {plot-data/mvo_p.txt};
  \end{axis}
\end{tikzpicture}
\end{center}
%
The portfolio has a well performing and smooth track record without any significant drawdowns over the period from 2013 to the middle of 2015.
The portfolio is managed for very risk averse investors that are willing to give up some expected return in order to maintain a very low level of volatility.

\begin{minipage}[t]{0.5\textwidth}
  \textbf{Portfolio Construction}
\begin{sit}\small
\item	Long-only equity exposure
\item	Fixed part invested in risk-free asset
\item	Yearly rebalanced
\item	Optimized for minimal risk
\item	Max. Allocation constraint
\item Hayashi-Yoshida Adjustment
\item Shrinkage of covariance estimates
\item Shrinkage of return estimates
\end{sit}
\end{minipage}
%
\begin{minipage}[t]{0.5\textwidth}
\textbf{Weights in Portfolio}\\
\tikzsetnextfilename{pf_10w}
\begin{tikzpicture}
  \begin{axis}[relw={0}{1}]
    \addplot table[x=time, y=value.DAX] {plot-data/mvo_relw.txt}\closedcycle;
    \addplot table[x=time, y=value.Dow.Jones] {plot-data/mvo_relw.txt}\closedcycle;
    \addplot table[x=time, y=value.Nikkei] {plot-data/mvo_relw.txt}\closedcycle;
    \addplot table[x=time, y=value.VXX] {plot-data/mvo_relw.txt}\closedcycle;
    \addplot table[x=time, y=value.mixed] {plot-data/mvo_relw.txt}\closedcycle;
  \end{axis}
\end{tikzpicture}
\end{minipage}

\begin{minipage}[t]{0.5\textwidth}
  \textbf{Opportunities}
\begin{sit}
    \item Prospect of positive return through participation in American, German and Japanese stock markets.
    \item Prospect of positive returns due to unhedged currency exposure.
    \item Reduced equity market exposure and protection to downturns due to risk free investments.
    \item Prospect of positive return through exploiting inefficiencies on equity option markets.
\end{sit}
\end{minipage}
\begin{minipage}[t]{0.5\textwidth}
  \textbf{Risks}
  \begin{sit}
    \item Losses possible in case stock markets develops against exposure built up in the portfolio.
    \item Losses possible if unhedged currency position moves against exposure built up in the portfolio.
    \item Underperformance due to risk free investments possible.
    \item Significant losses possible due to inefficiencies on equity option markets.% The volatility of the price of portfolio units might be increased.    
  \end{sit}
\end{minipage}

\newpage\section{JD Investments Equal Risk Contribution Fund}
\begin{mywraptable}{r}{0pt}
\begin{tabular}{lr}
\toprule
Ann. Return & 13.68\%\\
Ann. Volatility & 9.16\%\\
Sharpe Ratio & 1.49 \\
Max. Drawdown & 7.53\% \\
Turnover & 11.53\%\\
\bottomrule
\end{tabular}
\end{mywraptable}
\textit{JD Investments Equal Risk Contribution Fund} invests in the American, German and Japanese equity markets as well as the risk free rate.
The portfolio allocation aims to have each asset contributing equally to the portfolio risk.
The portfolio is rebalanced each month in order to be able to quickly adapt to changing a market environment.
%
This portfolio is particularly attractive to risk averse investors.
This portfolio completely ignores expected returns, making it robust for ever-changing markets or wrong assupmtions.
Investors can profit from higher than anticipated returns on less risky assets in the portfolio.
 

\tikzsetnextfilename{pf_6}
\begin{center}
\begin{tikzpicture}
  \begin{axis}[w={0}{170}, y filter/.expression={max(y,0)}]
    \addplot table[x=time, y=value.DAX] {plot-data/erc.txt}\closedcycle;
    \addlegendentry{DAX}
    \addplot table[x=time, y=value.Dow.Jones] {plot-data/erc.txt}\closedcycle;
    \addlegendentry{Dow Jones}
    \addplot table[x=time, y=value.Nikkei] {plot-data/erc.txt}\closedcycle;
    \addlegendentry{Nikkei}
    \pgfplotsset{cycle list shift=1}
    \addplot table[x=time, y=value.mixed] {plot-data/erc.txt}\closedcycle;
    \addlegendentry{Euribor}
    \addlegendimage{total} \addlegendentry{Total Value}
  \end{axis}
  \begin{axis}[w={0}{170}]
    \addplot[total] table[x=time, y=value] {plot-data/erc_p.txt};
  \end{axis}
\end{tikzpicture}
\end{center}

\begin{minipage}[t]{0.5\textwidth}
  \textbf{Portfolio Construction}
\begin{sit}\small
\item	Long-only equity exposure
\item	Investments in risk free assets
\item	Monthly rebalanced
\item	Risk parity optimization
\item Hayashi-Yoshida Adjustment
\item	Shrinkage of covariance estimates
\end{sit}
\end{minipage}
%
\begin{minipage}[t]{0.5\textwidth}
\textbf{Weights in Portfolio}\\
\tikzsetnextfilename{pf_6w}
\begin{tikzpicture}
  \begin{axis}[relw={0}{1}]
    \addplot table[x=time, y=value.DAX] {plot-data/erc_relw.txt}\closedcycle;
    %\addlegendentry{DAX}
    \addplot table[x=time, y=value.Dow.Jones] {plot-data/erc_relw.txt}\closedcycle;
    %\addlegendentry{Dow Jones}
    \addplot table[x=time, y=value.Nikkei] {plot-data/erc_relw.txt}\closedcycle;
    %\addlegendentry{Nikkei}
    \pgfplotsset{cycle list shift=1}
    \addplot table[x=time, y=value.mixed] {plot-data/erc_relw.txt}\closedcycle;
    %\addlegendentry{Euribor}
  \end{axis}
\end{tikzpicture}
\end{minipage}


The portfolio has smooth track record without any significant drawdowns over the period from 2013 to the middle of 2015.
It delivers an attractive return and low volatility.
The high level of diversification combined with the riskless asset promises the investor a stable and risk averse portfolio.
The portfolio is an attractive solution for risk averse investors looking for global, long only, equity exposure.
 

\begin{minipage}[t]{0.49\textwidth}
  \paragraph{Opportunities}
  \begin{sit}
    \item Prospect of positive return through participation in American, German and Japanese stock markets.
    \item Prospect of positive returns due to unhedged currency exposure.
    \item Reduced equity market exposure and protection to downturns due to risk free investments.
  \end{sit}
\end{minipage}
~~
\begin{minipage}[t]{0.49\textwidth}
  \paragraph{Risk}
  \begin{sit}
    \item Losses possible in case stock markets develop against exposure built up in the portfolio.
    \item Losses possible in case unhedged currency position develops against exposure built up in the portfolio.
    %\item Underperformance due to risk free investments possible.
    %\item Portfolio exposed to sudden jumps at the level of volatility.
    \item Sudden increase in maximum drawdowns of low risk stocks may cause considerable losses in the portfolio.
  \end{sit}
\end{minipage}

\newpage\section{JD Investments Market View Optimized Fund}
\begin{mywraptable}{r}{0pt}
\begin{tabular}{lr}
\toprule
Ann. Return & 17.71\%\\
Ann. Volatility & 12.00\%\\
Sharpe Ratio & 1.48\\
Max. Drawdown & 8.92\% \\
Turnover & 0.00\%\\
\bottomrule
\end{tabular}
\end{mywraptable}
\textit{JD Investments Market View Optimized Fund} invests in American, German and Japanese equity markets.
The fund seeks to minimize volatility while maximizing returns.
For this the average over the forward looking target prices of our assets was incorporated by using estimates made by analysts of 43 leading banks\footnote{In the beginning of 2013 the following expectations were used: German Market increases by 4.2\%, American Market 3.5\% and the Japanese Market by 2.00\%}.
The mix of fundamental analysis and systematic optimization makes this fund an attractive investment.


\tikzsetnextfilename{pf_5}
\begin{center}
\begin{tikzpicture}
  \begin{axis}[w={0}{170}, y filter/.expression={max(y,0)}]
    \addplot table[x=time, y=value.DAX] {plot-data/bl.txt}\closedcycle;
    \addlegendentry{DAX}
    \addplot table[x=time, y=value.Dow.Jones] {plot-data/bl.txt}\closedcycle;
    \addlegendentry{Dow Jones}
    \addplot table[x=time, y=value.Nikkei] {plot-data/bl.txt}\closedcycle;
    \addlegendentry{Nikkei}
    \addlegendimage{total} \addlegendentry{Total Value}
  \end{axis}
  \begin{axis}[w={0}{170}]
    \addplot[total] table[x=time, y=value] {plot-data/bl_p.txt};
  \end{axis}
\end{tikzpicture}
\end{center}

\begin{minipage}{0.5\textwidth}
\paragraph{Portfolio Construction}
\begin{sit}
\item Long only and full equity exposure
\item Fixed asset allocation
\item Optimized for Sharpe Ratio
\item Covariance estimator adjusted based on market views (using Black-Litterman)
\item No short-sales
\end{sit}
\end{minipage}
\begin{minipage}{0.5\textwidth}
\tikzsetnextfilename{pf_5w}
\begin{tikzpicture}
  \begin{axis}[relw={0}{1}, y filter/.expression={max(y,0)}]
    \addplot table[x=time, y=value.DAX] {plot-data/bl_relw.txt}\closedcycle;
    \addlegendentry{DAX}
    \addplot table[x=time, y=value.Dow.Jones] {plot-data/bl_relw.txt}\closedcycle;
    \addlegendentry{Dow Jones}
    \addplot table[x=time, y=value.Nikkei] {plot-data/bl_relw.txt}\closedcycle;
    \addlegendentry{Nikkei}
  \end{axis}
\end{tikzpicture}
\end{minipage}

The out of sample track record of the portfolio has is characterized by an attractive sharp ratio.
The fund is an appealing solution for moderate risk taking investors looking for global equity exposure.

\begin{minipage}[t]{0.49\textwidth}
  \paragraph{Opportunities}
  \begin{sit}
    \item Prospect of positive return through participation in American, German and Japanese stock markets.
    \item Prospect of positive returns due to unhedged currency exposure.
  \end{sit}
\end{minipage}
~~
\begin{minipage}[t]{0.49\textwidth}
  \paragraph{Risk}
  \begin{sit}
    \item Losses possible in case stock markets develop against exposure built up in the portfolio.
    \item Losses possible in case unhedged currency position develops against exposure built up in the portfolio.    
  \end{sit}
\end{minipage}

\newpage\section{JD Investments Low Risk Robust Fund}
\begin{mywraptable}{r}{0pt}
  \begin{tabular}{lr}
  \toprule
  Ann. Return & 18.10\%\\
  Ann. Volatility & 12.32\%\\
  Sharpe & 1.47\\
  Max. Drawdown& 8.20\%\\
  Turnover & 9.47\%\\
  \bottomrule
  \end{tabular}
\end{mywraptable}
\textit{JD Investments Low Risk Robust} Fund has full exposure to the American, German and Japanese equity markets.
The asset allocation is aimed at minimizing volatility while variances/covariance matrix is estimated using robust covariance-shrinkage method.
This portfolio still provides attractive returns and Sharpe ratio in spite of being oriented on variance minimization.

\tikzsetnextfilename{pf_7}
\begin{center}
\begin{tikzpicture}
  \begin{axis}[w={0}{170}, y filter/.expression={max(y,0)}]
    \addplot table[x=time, y=value.DAX] {plot-data/minv_rob.txt}\closedcycle;
    \addlegendentry{DAX}
    \addplot table[x=time, y=value.Dow.Jones] {plot-data/minv_rob.txt}\closedcycle;
    \addlegendentry{Dow Jones}
    \addplot table[x=time, y=value.Nikkei] {plot-data/minv_rob.txt}\closedcycle;
    \addlegendentry{Nikkei}
    \addlegendimage{total} \addlegendentry{Total Value}
  \end{axis}
  \begin{axis}[w={0}{170}]
    \addplot[total] table[x=time, y=value] {plot-data/minv_rob_p.txt};
  \end{axis}
\end{tikzpicture}
\end{center}

\begin{minipage}{0.5\textwidth}
\paragraph{Portfolio Construction}
\begin{sit}
\item Long only and full equity exposure
\item Yearly rebalanced
\item Robust estimate of covariance
\item Optimized for Minimum Risk
\item Robust for different market environments
\item No short-sales
\end{sit}
\end{minipage}
\tikzsetnextfilename{pf_7w}
\begin{minipage}{0.5\textwidth}
\begin{tikzpicture}
\begin{axis}[relw={0}{1}]
  \addplot table[x=time, y=value.DAX] {plot-data/minv_rob_relw.txt}\closedcycle;
  \addlegendentry{DAX}
  \addplot table[x=time, y=value.Dow.Jones] {plot-data/minv_rob_relw.txt}\closedcycle;
  \addlegendentry{Dow Jones}
  \addplot table[x=time, y=value.Nikkei] {plot-data/minv_rob_relw.txt}\closedcycle;
  \addlegendentry{Nikkei}
\end{axis}
\end{tikzpicture}
\end{minipage}


The portfolio has a well performing and smooth track record without any significant drawdowns over the period from 2013 to the middle of 2015.
An attractive return with moderate volatility is the key characteristic of this fund.
 

\begin{minipage}[t]{0.49\textwidth}
  \paragraph{Opportunities}
  \begin{sit}
  \item Prospect of positive return through participation in American, German and Japanese stock markets.
  \item Prospect of positive returns due to unhedged currency exposure. Stocks may cause considerable losses in the portfolio.
  \end{sit}
\end{minipage}
~~
\begin{minipage}[t]{0.49\textwidth}
  \paragraph{Risk}
  \begin{sit}
  \item Losses possible in case stock markets develops against exposure built up in the portfolio.
  \item Losses possible if unhedged currency position moves against exposure built up in the portfolio.
  \end{sit}
\end{minipage}

\newpage\section{JD Investments Fixed Weight Equity Fund}
\begin{mywraptable}{r}{0pt}
\begin{tabular}{lr}
\toprule
Ann. Return & 17.96\%\\
Ann. Volatility & 12.61\%\\
Sharpe Ratio & 1.42 \\
Max. Drawdown & 9.08\%\\
Turnover & 0.00\%\\
\bottomrule
\end{tabular}
\end{mywraptable}
\textit{JD Investments Fixed Weight Equity Fund} invests in American, German and Japanese equity markets. The asset allocation is static and most diversified while disregarding expected returns and variances/covariances. This portfolio is especially appealing since it is resistant against errors in true variance/covariance and expected return estimation.


\begin{center}
\tikzsetnextfilename{pf_3}
\begin{tikzpicture}
  \begin{axis}[w={0}{170}]
    \addplot table[x=time, y=value.DAX] {plot-data/fw.txt}\closedcycle;
    \addlegendentry{DAX}
    \addplot table[x=time, y=value.Dow.Jones] {plot-data/fw.txt}\closedcycle;
    \addlegendentry{Dow Jones}
    \addplot table[x=time, y=value.Nikkei] {plot-data/fw.txt}\closedcycle;
    \addlegendentry{Nikkei}
    \addlegendimage{total} \addlegendentry{Total Value}
  \end{axis}
  \begin{axis}[w={0}{170}]
    \addplot[total] table[x=time, y=value] {plot-data/fw_p.txt};
  \end{axis}
\end{tikzpicture}
\end{center}

\begin{minipage}{0.5\textwidth}
\paragraph{Portfolio Construction}
\begin{sit}
\item	Long only and full equity exposure
\item	Fixed asset allocation
\item	Strong diversification
\item	No assumption on expected returns and variances/covariances
\end{sit}
\end{minipage}
\begin{minipage}{0.5\textwidth}
\tikzsetnextfilename{pf_3w}
\begin{tikzpicture}
  \begin{axis}[relw={0}{1}]
    \addplot table[x=time, y=value.DAX] {plot-data/fw_relw.txt}\closedcycle;
    \addlegendentry{DAX}
    \addplot table[x=time, y=value.Dow.Jones] {plot-data/fw_relw.txt}\closedcycle;
    \addlegendentry{Dow Jones}
    \addplot table[x=time, y=value.Nikkei] {plot-data/fw_relw.txt}\closedcycle;
    \addlegendentry{Nikkei}
  \end{axis}
\end{tikzpicture}
\end{minipage}


With a promising track record and a solid risk return profile as well as zero turnover and therefore very appealing net cost return the \textit{JD Investments Fixed Weight Equity Fund} is an attractive investment for the investor willing to take on higher risks in exchange for higher returns. 

\begin{minipage}[t]{0.49\textwidth}
  \paragraph{Opportunities}
  \begin{sit}
  \item Prospect of positive return through participation in American, German and Japanese stock markets.
  \item Prospect of positive returns due to unhedged currency exposure. Stocks may cause considerable losses in the portfolio.
  \end{sit}
\end{minipage}
~~
\begin{minipage}[t]{0.49\textwidth}
  \paragraph{Risk}
  \begin{sit}
  \item Losses possible in case stock markets develops against exposure built up in the portfolio.
  \item Losses possible if unhedged currency position moves against exposure built up in the portfolio.
  \end{sit}
\end{minipage}

\setlength{\intextsep}{\oldintextsep}
\chapter{Conclusion \& Recommendation}
To make an in-depth analysis of the performance of selected portfolios we take a look at their out-of-sample performance.

\begin{table}[ht]
\centering
\begin{tabular}{lrrrrrrr}
\toprule
                      & Sharpe & $\mu$[\%] &$\sigma$[\%]& MDD[\%]& T/O[\%] & TC[\%] & $\alpha$[\%] \\ 
\midrule
 Mean Variance+ & 1.39 & 11.61 & 8.35 & 5.88 & 10.34 & 0.05 & 0.38 \\
 Equal Risk Contribution& 1.49 & 13.68 & 9.16 & 7.53 & 11.53 & 0.06 & 0.64 \\ 
 Market View Optimized & 1.48 & 17.71 & 12.00 & 8.92 & 0.00 & 0.00 & 0.55 \\ 
 Low Risk Robust & 1.47 & 18.10 & 12.32 & 8.20 & 9.47 & 0.05 & 0.44 \\ 
 Fixed Weights & 1.42 & 17.96 & 12.61 & 9.08 & 0.00 & 0.00 & 0.62 \\ 
%Minimum Variance & 1.37 & 10.92 & 7.99 & 5.29 & 11.56 & 0.06 & 0.28 \\ 
 %Maximum Sharpe & 1.46 & 17.55 & 12.05 & 8.98 & 11.94 & 0.06 & 0.26 \\ 
 %Fixed Weights Leveraged & 1.54 & 21.47 & 13.95 & 9.83 & 0.00 & 0.00 & 0.76 \\ 
 %Maximum Sharpe Robust & 1.12 & 16.55 & 14.82 & 7.70 & 44.45 & 0.22 & 0.26 \\ 
 %MV Max Allocation & 1.40 & 5.18 & 3.70 & 2.50 & 4.65 & 0.02 & 0.19 \\ 
 \bottomrule
\end{tabular}
\end{table}

\begin{figure}
\tikzsetnextfilename{comparison}
\begin{tikzpicture}
  \begin{axis}[w={0}{180}, stack plots = false,xmin={2013-01-01}, xmax={2015-06-30}, ymin = 95, legend columns = 1]
    \addplot[total, red] table[x=time, y=value] {plot-data/fw_p.txt};
    \addlegendentry{Fixed Weights}
    \addplot[total, blue] table[x=time, y=value] {plot-data/bl_p.txt};
    \addlegendentry{Market View Optimized}
    \addplot[total, green] table[x=time, y=value] {plot-data/erc_p.txt};
    \addlegendentry{Equal Risk Contribution}
    \addplot[total, orange] table[x=time, y=value] {plot-data/minv_rob_p.txt};
    \addlegendentry{Low Risk Robust}
    \addplot[total, black] table[x=time, y=value] {plot-data/mvo_p.txt};
    \addlegendentry{Mean Variance+}
    \addplot[total, gray] table[x=time, y=value] {plot-data/msci.txt};
    \addlegendentry{MSCI}
  \end{axis}
\end{tikzpicture}
\end{figure}


Compared to the \textit{JD Investments Black Litterman Equity Fund} our \textit{JD Investments Fixed Weight Equity Fund} has with 17.96\% a higher rate of return. However, its sharpe ratio is lower. The maximum drawdown is also higher lying at 9.08\%. Both have a turnover of 0\% due to their static asset allocation.

The \textit{JD Investments Equal Risk Contribution Equity Fund} has with 13.68\% a lower annual return than our \textit{JD Investments Black Litterman Equity Fund} and \textit{JD Investments Fixed Weight Equity Fund}. It also has a lower volatility and maximum drawdown. With 1.49 it has a higher sharpe ratio than the latter two. Turnover is 11.53\% with transaction costs of 0.06\%. 

Of the five portfolios introduced the \textit{JD Investments Robust Minimum Variance Equity Fund} has the highest attainable performance. The \textit{JD Investments Robust Minimum Variance Equity Fund} has a higher sharpe ratio than the \textit{JD Investments Fixed Weight Equity Fund} but it is lower than for the \textit{JD Investments Equal Risk Contribution Equity Fund} and \textit{JD Investments Black Litterman Equity Fund}. The maximum drawdown is lower than for \textit{JD Investments Black Litterman Equity Fund} and \textit{JD Investments Fixed Weight Equity Fund} but higher than the equal risk portfolio. The turnover is higher than for \textit{JD Investments Black Litterman Equity Fund} and \textit{JD Investments Fixed Weight Equity Fund} but lower than the \textit{JD Investments Robust Minimum Variance Equity Fund}. Transaction costs amount to 0.05\%.

The \textit{JD Investments Mean Variance+ Equity Fund} has the lowest return compared to the other four funds and the lowest sharpe ratio. Respectively, it also has the lowest volatility and maximum drawdown. Turnover is 12.10\% and transaction costs are 0.06\%. Given the discussion above, it can be seen that choosing an optimal portfolio is not a straightforward task as it might look on first sight. The decision for the optimal portfolio is based on investor’s constraints and implied risk aversion. 

After comparing the performance metrics of the five portfolios we believe that the \textit{JD Investments Equal Risk Contribution Equity Fund} is the best portfolio choice for you. Even though it has a lower performance than the \textit{JD Investments Fixed Weight Equity Fund}, \textit{JD Investments Black Litterman Equity Fund} and \textit{JD Investments Robust Minimum Variance Equity Fund} we believe it fits best to your investment needs. Of the five portfolios it has the highest sharpe ratio which means that it gives you the highest risk-return trade-off. Secondly, it is the only portfolio that generates a statistically significant positive alpha. This means that it is evident that you are benefiting from us as we proved to make solid investment decisions.

\chapter{Bonus: 2015-06-30 --- 2015-11-01}
\begin{table}[ht]
\centering
\begin{tabular}{lrrrrrr}
  \toprule
 & Sharpe & $\mu$[\%] & $\sigma$[\%]& Inf. Ratio & Appraisal Ratio\\
  \midrule
   Mean Variance+ & -1.03 & -3.93 & 3.80 & -1.08 & 0.16\\ 
   Equal Risk Contribution & -0.81 & -11.30 & 13.90 & -0.81 & 0.14\\ 
   Market View Optimized & -1.20 & -22.38 & 18.64 & -2.11 & 0.11\\ 
   Low Risk Robust& -1.38 & -25.32 & 18.34 & -2.24 & 0.12\\ 
   Fixed Weights & -1.20 & -22.38 & 18.64 & -2.04 & 0.12\\ 
%Minimum Variance & -1.32 & -15.51 & 11.76 & 0.93\\ 
  %Maximum Sharpe & -1.14 & -21.09 & 18.51 & -0.42\\ 
  %Fixed Weights Leveraged & -1.27 & -25.50 & 20.02 & -0.94\\ 
  %Max Sharpe Robust & -1.21 & -22.29 & 18.38 & -0.58\\ 
  %MV Max Allocation & -1.03 & -3.87 & 3.74 & 1.52\\ 
   \bottomrule
\end{tabular}
\end{table}

\tikzsetnextfilename{comparison_bonus}
\begin{tikzpicture}
  \begin{axis}[w={80}{110}, stack plots = false, xtick={2015-07-01, 2015-08-01, 2015-09-01, 2015-10-01, 2015-11-01}, xticklabel={\pgfcalendar{tickcal}{\tick}{\tick}{\pgfcalendarshorthand{m}{.}}}, legend columns = 1]
    \addplot[total, red] table[x=time, y=value.Fixed.Weights] {plot-data/comparison_bonus.txt};
    \addlegendentry{Fixed Weights}
    \addplot[total, blue] table[x=time, y=value.Black.Litterman] {plot-data/comparison_bonus.txt};
    \addlegendentry{Market View Optimized}
    \addplot[total, green] table[x=time, y=value.Equal.Risk] {plot-data/comparison_bonus.txt};
    \addlegendentry{Equal Risk Contribution}
    \addplot[total, orange] table[x=time, y=value.Min.Variance.Robust] {plot-data/comparison_bonus.txt};
    \addlegendentry{Low Risk Robust}
    \addplot[total, black] table[x=time, y= value.Mean.Variance..Equity.Fund] {plot-data/comparison_bonus.txt};
    \addlegendentry{Mean Variance+}
    \addplot[total, gray] table[x=time, y=value.msci] {plot-data/comparison_bonus.txt};
    \addlegendentry{MSCI}
  \end{axis}
\end{tikzpicture}




%\chapter{Out Of Sample: 2013-01-01 --- 2015-06-30}
%\begin{table}[ht]
%\centering
%\begin{tabular}{lrrrrrrr}
%\toprule
                      %& Sharpe & $\mu$ &$\sigma$& MDD & T/O & TC & $\alpha$ \\ 
%\midrule
%Minimum Variance & 1.37 & 10.92 & 7.99 & 5.29 & 11.56 & 0.06 & -0.56 \\ 
  %Maximum Sharpe & 1.46 & 17.55 & 12.05 & 8.98 & 11.94 & 0.06 & -0.65 \\ 
  %Fixed Weights & 1.42 & 17.96 & 12.61 & 9.08 & 0.00 & 0.00 &  \\ 
  %Fixed Weights Leveraged & 1.54 & 21.47 & 13.95 & 9.83 & 0.00 & 0.00 &  \\ 
  %Black Litterman & 1.48 & 17.71 & 12.00 & 8.92 & 0.00 & 0.00 &  \\ 
  %Equal Risk & 1.49 & 13.68 & 9.16 & 7.53 & 11.53 & 0.06 & 0.33 \\ 
  %Min Variance Robust & 1.47 & 18.10 & 12.32 & 8.20 & 9.47 & 0.05 & -0.66 \\ 
  %Max Sharpe Robust & 1.12 & 16.55 & 14.82 & 7.70 & 44.45 & 0.22 & -0.49 \\ 
  %MV Max Allocation & 1.31 & 4.12 & 3.15 & 3.05 & 10.41 & 0.05 & -0.45 \\ 
  %MV Max Allocation \& Short-Constraint & 1.40 & 12.34 & 8.79 & 6.09 & 12.10 & 0.06 & -0.67 \\ 
%\bottomrule
%\end{tabular}
%\end{table}

\newpage
\chapter{Appendix}
\section{1 -- JD Investments Minimum Variance Equity Fund}
\textit{JD Investments Minimum Variance Equity} fund invests in American, German and Japanese equity market together with 6M Euribor to substantially reduce risk.
The asset allocation is oriented on volatility minimization while disregarding expected returns.
This portfolio is particularly suitable for the clients with strong risk-aversion, since it does not try to speculate on expected returns but rather cares just about reducing variance.

\paragraph{Portfolio Construction}
\begin{sit}
\item Long only equity exposure

\item High concentration in risk-free investment (45\% of portfolio)
\item No assumption on expected returns
\end{sit}

\tikzsetnextfilename{pf_1}
\begin{figure}[H]
\begin{tikzpicture}
  \begin{axis}[w={0}{140}, y filter/.expression={max(y,0)}]
    \addplot table[x=time, y=value.DAX] {plot-data/minv.txt}\closedcycle;
    \addlegendentry{DAX}
    \addplot table[x=time, y=value.Dow.Jones] {plot-data/minv.txt}\closedcycle;
    \addlegendentry{Dow Jones}
    \addplot table[x=time, y=value.Nikkei] {plot-data/minv.txt}\closedcycle;
    \addlegendentry{Nikkei}
    \pgfplotsset{cycle list shift=1}
    \addplot table[x=time, y=value.mixed] {plot-data/minv.txt}\closedcycle;
    \addlegendentry{Euribor}
    \addlegendimage{total} \addlegendentry{Total Value}
  \end{axis}
  \begin{axis}[w={0}{140}, stack plots = false]
    \addplot[total] table[x=time, y=value] {plot-data/minv_p.txt};
    %\addplot[black] table[x=time, y=value] {plot-data/msci.txt};
  \end{axis}
  \begin{axis}[w={0}{140}, y filter/.expression={min(y,0)}]
    \addplot table[x=time, y=value.DAX] {plot-data/minv.txt}\closedcycle;
    \addplot table[x=time, y=value.Dow.Jones] {plot-data/minv.txt}\closedcycle;
    \addplot table[x=time, y=value.Nikkei] {plot-data/minv.txt}\closedcycle;
    \pgfplotsset{cycle list shift=4}
    \addplot table[x=time, y=value.mixed] {plot-data/minv.txt}\closedcycle;
  \end{axis}
\end{tikzpicture}
\end{figure}

\tikzsetnextfilename{pf_1w}
\vspace{\abovedisplayskip}
\begin{minipage}{0.65\textwidth}
\begin{tikzpicture}
  \begin{axis}[relw={0}{1}, y filter/.expression={max(y,0)}]
    \addplot table[x=time, y=value.DAX] {plot-data/minv_relw.txt}\closedcycle;
    \addlegendentry{DAX}
    \addplot table[x=time, y=value.Dow.Jones] {plot-data/minv_relw.txt}\closedcycle;
    \addlegendentry{Dow Jones}
    \addplot table[x=time, y=value.Nikkei] {plot-data/minv_relw.txt}\closedcycle;
    \addlegendentry{Nikkei}
    \pgfplotsset{cycle list shift=1}
    \addplot table[x=time, y=value.mixed] {plot-data/minv_relw.txt}\closedcycle;
    \addlegendentry{Euribor}
  \end{axis}
  \begin{axis}[relw={0}{1},  y filter/.expression={min(y,0)}]
    \addplot table[x=time, y=value.DAX] {plot-data/minv_relw.txt}\closedcycle;
    \addplot table[x=time, y=value.Dow.Jones] {plot-data/minv_relw.txt}\closedcycle;
    \addplot table[x=time, y=value.Nikkei] {plot-data/minv_relw.txt}\closedcycle;
    \pgfplotsset{cycle list shift=1}
    \addplot table[x=time, y=value.mixed] {plot-data/minv_relw.txt}\closedcycle;
  \end{axis}
\end{tikzpicture}
\end{minipage}
\begin{minipage}{0.35\textwidth}
\begin{tabular}{lr}
\toprule
Ann. Return & 10.92\%\\
Ann. Volatility & 7.99\%\\
Sharpe Ratio & 1.37\\
Max. Drawdown & 5.92\% \\
Turnover & 11.56\%\\
\bottomrule
\end{tabular}
\end{minipage}

This portfolio is suitable for very risk-averse clients who are ready to sacrifice potential returns and have lower Sharpe ratio just to reduce volatility.

\begin{table}[H]
\begin{tabularx}{\textwidth}{XX}
  \toprule
  \textbf{\textsf{Opportunities}} & \textbf{\textsf{Risks}} \\
  \midrule
Possibility to gain attractive returns through participation American, German and Japanese stock markets.&
Risks of obtaining very low returns in the periods of high concentration in Euribor since  return on short-term Euribor is often 0 and sometimes even negative. \\[1em]
Risks of indexes are substantially reduced with the help of mixture with risk-free asset. &
Losses possible in case unhedged currency position develop against exposure built up in the portfolio.\\
  \bottomrule
\end{tabularx}
\end{table}


\newpage\section{2 -- Max Sharpe}
The Maximum Sharpe Ratio portfolio is constructed with the use of sophisticated mathematical techniques. The objective of the fund is to deliver significant risk adjusted return, and to outperform any other portfolio within a similar investment universe. 
At any point in time, assets of the portfolio are invested fully into risky assets. It is suitable for enthusiastic risk seekers than are willing to take high risk to achieve significant returns. As shown in the figure below, majority of funds are invested into US market and the rest is allocated in Germany and Japan. From the beginning of 2013 to the middle of 2015, MSR portfolio delivered return of over 17\% annually, while volatility was moderate and amounted to 12.05\%. 

\tikzsetnextfilename{pf_2}
\begin{figure}[H]
\begin{tikzpicture}
  \begin{axis}[w={0}{170}, y filter/.expression={max(y,0)}]
    \addplot table[x=time, y=value.DAX] {plot-data/ms.txt}\closedcycle;
    \addlegendentry{DAX}
    \addplot table[x=time, y=value.Dow.Jones] {plot-data/ms.txt}\closedcycle;
    \addlegendentry{Dow Jones}
    \addplot table[x=time, y=value.Nikkei] {plot-data/ms.txt}\closedcycle;
    \addlegendentry{Nikkei}
    \addlegendimage{total} \addlegendentry{Total Value}
  \end{axis}
  \begin{axis}[w={0}{170}]
    \addplot[total] table[x=time, y=value] {plot-data/ms_p.txt};
  \end{axis}
  \begin{axis}[w={0}{170}, y filter/.expression={min(y,0)}]
    \addplot table[x=time, y=value.DAX] {plot-data/ms.txt}\closedcycle;
    \addplot table[x=time, y=value.Dow.Jones] {plot-data/ms.txt}\closedcycle;
    \addplot table[x=time, y=value.Nikkei] {plot-data/ms.txt}\closedcycle;
  \end{axis}
\end{tikzpicture}
\end{figure}

\tikzsetnextfilename{pf_2w}
\vspace{\abovedisplayskip}
\begin{minipage}{0.65\textwidth}
  \begin{tikzpicture}
  \begin{axis}[relw={0}{1}, y filter/.expression={max(y,0)}]
    \addplot table[x=time, y=value.DAX] {plot-data/ms_relw.txt}\closedcycle;
    \addlegendentry{DAX}
    \addplot table[x=time, y=value.Dow.Jones] {plot-data/ms_relw.txt}\closedcycle;
    \addlegendentry{Dow Jones}
    \addplot table[x=time, y=value.Nikkei] {plot-data/ms_relw.txt}\closedcycle;
    \addlegendentry{Nikkei}
  \end{axis}
  \begin{axis}[relw={0}{1}, y filter/.expression={min(y,0)}]
    \addplot table[x=time, y=value.DAX] {plot-data/ms_relw.txt}\closedcycle;
    \addplot table[x=time, y=value.Dow.Jones] {plot-data/ms_relw.txt}\closedcycle;
    \addplot table[x=time, y=value.Nikkei] {plot-data/ms_relw.txt}\closedcycle;
  \end{axis}
\end{tikzpicture}
\end{minipage}
\begin{minipage}{0.35\textwidth}
\begin{tabular}{lr}
\toprule
Ann. Return & 17.55\%\\
Ann. Volatility & 12.05\%\\
Sharpe Ratio & 1.46 \\
Max. Drawdown & 8.92\% \\
Turnover & 11.94\%\\
\bottomrule
\end{tabular}
\end{minipage}



\begin{table}[H]
\begin{tabularx}{\textwidth}{XX}
  \toprule
  \textbf{\textsf{Opportunities}} & \textbf{\textsf{Risks}} \\
  \midrule
Active participation in the growth of US, German and Japanese stock market. Investor profits from momentum in the stock market. &
Portfolio exposed to sudden increase in volatility. \\[1em]
Portfolio rebalanced and optimized on yearly basis. By optimizing, investor achieves “best” asset allocation given investment set. &
Due to 100\% allocation in risky assets, investor exposed to liquidity shock in the market.\\
  \bottomrule
\end{tabularx}
\end{table}

\newpage\section{4 -- Equal Weight Leveraged}
EWL portfolio offers investor asset allocation into three different markets, namely US, German and Japanese. Weights in the portfolio are equally distributed with each asset class consisting of 36\% of total portfolio value. Furthermore, investor’s return will be enhanced by shorting VXX index that is predicted to have a bear run in the near future. Portfolio has no holdings in riskless assets, and invests heavily into diversified market indices. As shown in the figure below, portfolio has performed extremely well over the past two years. In this period EWL portfolio achieved return of more than 21\% with considerably lower volatility. 

\tikzsetnextfilename{pf_4}
\begin{figure}[H]
\begin{tikzpicture}
  \begin{axis}[w={-20}{190}, y filter/.expression={max(y,0)}]
    \addplot table[x=time, y=value.DAX] {plot-data/fw_lev.txt}\closedcycle;
    \addlegendentry{DAX}
    \addplot table[x=time, y=value.Dow.Jones] {plot-data/fw_lev.txt}\closedcycle;
    \addlegendentry{Dow Jones}
    \addplot table[x=time, y=value.Nikkei] {plot-data/fw_lev.txt}\closedcycle;
    \addlegendentry{Nikkei}
    \addplot table[x=time, y=value.VXX] {plot-data/fw_lev.txt}\closedcycle;
    \addlegendentry{VXX}
    \addlegendimage{total} \addlegendentry{Total Value}
  \end{axis}
  \begin{axis}[w={-20}{190}]
    \addplot[total] table[x=time, y=value] {plot-data/fw_lev_p.txt};
  \end{axis}
  \begin{axis}[w={-20}{190}, bottom, y filter/.expression={min(y,0)}]
    \pgfplotsset{cycle list shift=3}
    \addplot table[x=time, y=value.VXX] {plot-data/fw_lev.txt}\closedcycle;
  \end{axis}
\end{tikzpicture}
\end{figure}

\tikzsetnextfilename{pf_4w}
\vspace{\abovedisplayskip}
\begin{minipage}{0.65\textwidth}
\begin{tikzpicture}
  \begin{axis}[relw={-0.1}{1.2}, y filter/.expression={max(y,0)}]
    \addplot table[x=time, y=value.DAX] {plot-data/fw_lev_relw.txt}\closedcycle;
    \addlegendentry{DAX}
    \addplot table[x=time, y=value.Dow.Jones] {plot-data/fw_lev_relw.txt}\closedcycle;
    \addlegendentry{Dow Jones}
    \addplot table[x=time, y=value.Nikkei] {plot-data/fw_lev_relw.txt}\closedcycle;
    \addlegendentry{Nikkei}
    \addplot table[x=time, y=value.VXX] {plot-data/fw_lev_relw.txt}\closedcycle;
    \addlegendentry{VXX}
  \end{axis}
  \begin{axis}[relw={-0.1}{1.2}, bottom, y filter/.expression={min(y,0)}]
    \pgfplotsset{cycle list shift=3}
    \addplot table[x=time, y=value.VXX] {plot-data/fw_lev_relw.txt}\closedcycle;
  \end{axis}
\end{tikzpicture}
\end{minipage}
\begin{minipage}{0.35\textwidth}
\begin{tabular}{lr}
\toprule
Ann. Return & 21.47\%\\
Ann. Volatility & 13.95\%\\
Sharpe Ratio & 1.54 \\
Max. Drawdown & 9.83\% \\
Turnover & 0\\
\bottomrule
\end{tabular}
\end{minipage}

\begin{table}[H]
\begin{tabularx}{\textwidth}{XX}
  \toprule
  \textbf{\textsf{Opportunities}} & \textbf{\textsf{Risks}} \\
  \midrule
Active participation in the growth of US, German and Japanese stock market. Investor profits low volatility in the stock market. &
Portfolio is extremely exposed to sudden increase in volatility, since it is short volatility index.\\[1em]
Portfolio is not rebalanced and thus there are no extra fees for buying or selling securities. &
Maximal drawdown has to be carefully monitored due to the running of short position.\\
  \bottomrule
\end{tabularx}
\end{table}

%\newpage\section{8 --- Max Sharpe Robust}
%\begin{tikzpicture}
  %\begin{axis}[w={0}{170}, y filter/.expression={max(y,0)}]
    %\addplot table[x=time, y=value.DAX] {plot-data/msr_rob.txt}\closedcycle;
    %\addlegendentry{DAX}
    %\addplot table[x=time, y=value.Dow.Jones] {plot-data/msr_rob.txt}\closedcycle;
    %\addlegendentry{Dow Jones}
    %\addplot table[x=time, y=value.Nikkei] {plot-data/msr_rob.txt}\closedcycle;
    %\addlegendentry{Nikkei}
    %\addlegendimage{total} \addlegendentry{Total Value}
  %\end{axis}
  %\begin{axis}[w={0}{170}]
    %\addplot[total] table[x=time, y=value] {plot-data/msr_rob_p.txt};
  %\end{axis}
%\end{tikzpicture}

%\begin{tikzpicture}
  %\begin{axis}[relw={0}{1}]
    %\addplot table[x=time, y=value.DAX] {plot-data/msr_rob_relw.txt}\closedcycle;
    %\addlegendentry{DAX}
    %\addplot table[x=time, y=value.Dow.Jones] {plot-data/msr_rob_relw.txt}\closedcycle;
    %\addlegendentry{Dow Jones}
    %\addplot table[x=time, y=value.Nikkei] {plot-data/msr_rob_relw.txt}\closedcycle;
    %\addlegendentry{Nikkei}
  %\end{axis}
%\end{tikzpicture}

\newpage\section{9 --- Mean Variance with Maximum Allocation Constraint}
\begin{figure}[H]
\begin{tikzpicture}
  \begin{axis}[w={-10}{120}, y filter/.expression={max(y,0)}]
    \addplot table[x=time, y=value.DAX] {plot-data/mva.txt}\closedcycle;
    \addlegendentry{DAX}
    \addplot table[x=time, y=value.Dow.Jones] {plot-data/mva.txt}\closedcycle;
    \addlegendentry{Dow Jones}
    \addplot table[x=time, y=value.Nikkei] {plot-data/mva.txt}\closedcycle;
    \addlegendentry{Nikkei}
    \addplot table[x=time, y=value.VXX] {plot-data/mva.txt}\closedcycle;
    \addlegendentry{VXX}
    \addplot table[x=time, y=value.mixed] {plot-data/mva.txt}\closedcycle;
    \addlegendentry{Euribor}
    \addlegendimage{total} \addlegendentry{Total Value}
  \end{axis}
  \begin{axis}[w={-10}{120}]
    \addplot[total] table[x=time, y=value] {plot-data/mva_p.txt};
  \end{axis}
  \begin{axis}[w={-10}{120}, bottom, y filter/.expression={min(y,0)}]
    \pgfplotsset{cycle list shift=3}
    \addplot table[x=time, y=value.VXX] {plot-data/mva.txt}\closedcycle;
  \end{axis}
\end{tikzpicture}
\end{figure}

\begin{tikzpicture}
  \begin{axis}[relw={-0.1}{1.1}, y filter/.expression={max(y,0)}]
    \addplot table[x=time, y=value.DAX] {plot-data/mva_relw.txt}\closedcycle;
    \addlegendentry{DAX}
    \addplot table[x=time, y=value.Dow.Jones] {plot-data/mva_relw.txt}\closedcycle;
    \addlegendentry{Dow Jones}
    \addplot table[x=time, y=value.Nikkei] {plot-data/mva_relw.txt}\closedcycle;
    \addlegendentry{Nikkei}
    \addplot table[x=time, y=value.VXX] {plot-data/mva_relw.txt}\closedcycle;
    \addlegendentry{VXX}
    \addplot table[x=time, y=value.mixed] {plot-data/mva_relw.txt}\closedcycle;
    \addlegendentry{Euribor}
  \end{axis}
  \begin{axis}[relw={-0.1}{1.1}, bottom, y filter/.expression={min(y,0)}]
    \pgfplotsset{cycle list shift=3}
    \addplot table[x=time, y=value.VXX] {plot-data/mva_relw.txt}\closedcycle;
  \end{axis}
\end{tikzpicture}


\chapter{Code}
We developed our asset allocation algorithms in R. In the following section, we would like to briefly explain the basic rationaly behind each module of our software. Each module fulfills a disinct function:
\begin{description}
  \item[0-helper] Support of data visualization and export
  \item[1-save data] Download of historic data from Yahoo Finance, Quandl and other sources
  \item[1-data] Basic data preparation tasks
  \item[2-analysis] Model specification, regularization techniques, Evaluation logic
  \item[2-application] Composition of portfolios through abstraction of portfolio models and backtesting parameters into a ``job'' data structure; Data structure is passed into a ``run job'' function which returns a data structure containing, (a) individual asset weights, (b) overall portfolio value (c) generated portfolio weights (d) computed performance measures (e) original job specification.\\
This data structure can subsequently be processed/evaluated appropriately. The structure of modules and functions as well as the abstraction into jobs allows for a clear separation of model specification and application on historic data. The historical simulation of a portfolio reduces to the specification of the model and backtesting setting by means of composition of idempotent functions.
  \item[3-analysis2] Numerical optimization algorithm to incorporate maximum allocation/short-sale constraints
  \item[mdd] Algorithm for calculating the maximum drawdown of a given timeseries
\end{description}

\lstinputlisting[language=R, basicstyle=\scriptsize]{../0-helper.R}

\lstinputlisting[language=R, basicstyle=\scriptsize]{../1-save_data.R}

\lstinputlisting[language=R, basicstyle=\scriptsize]{../1-data.R}

\lstinputlisting[language=R, basicstyle=\scriptsize]{../2-analysis.R}

\lstinputlisting[language=R, basicstyle=\scriptsize]{../2-application.R}

\lstinputlisting[language=R, basicstyle=\scriptsize]{../3-analysis2.R}

\lstinputlisting[language=R, basicstyle=\scriptsize]{../mdd.R}

\end{document}
